\chapter{تطبيقات الغربال والفجوات بين الأعداد الأولية}
\symbolindex{تطبيقات الغربال والفجوات بين الأعداد الأولية}

\section{نطاق الفصل وحالته}

\noindent\textbf{حالة الفصل:}
\textenglish{REVIEWED / NOT-RELEASE-READY}.
أُنشئ هذا المتن بعد إغلاق بوابة ما قبل التأليف وصدور
\textenglish{PASS-FOR-AUTHORING = YES}. اكتمل تدقيق المصادر الأولية وعدم الدور،
وأُدخلت مفاتيح المراجع الأساسية، واجتاز الفصل تدقيق ما بعد التأليف والبناء الكامل
(\textenglish{XeLaTeX -> Biber -> XeLaTeX -> XeLaTeX}، بلا إحالات أو مراجع غير معرّفة).

أعادت المراجعة العلمية المستقلة فحص المتن كاملًا وأصدرت الحكم
\textenglish{APPROVED} مع ملاحظات طباعية غير حاجزة، إذ لم تظهر عوائق رياضية أو
مرجعية أو حوكمية أو بنائية. راجع المالك المتن على هذا الأساس واعتمده، فيحمل
الفصل حالة \textenglish{REVIEWED}، لكنه لا يحمل بعد حالة
\textenglish{RELEASE-READY}؛ يبقى الدمج في \texttt{main} قرارًا منفصلاً للمالك.

ينتقل هذا الفصل من الحد العلوي الغربالي المجرد في الفصل الخامس عشر إلى تطبيقات
على الأزواج الأولية والأعداد شبه الأولية والفجوات بين الأوليات. ويُفصل صراحة بين
ما يُثبت هنا وما يُقتبس من الأدبيات الأصلية.

لا يثبت الفصل حدسية الأوليات التوأم، ولا يثبت \(H_1=2\)، ولا يعرض ثابتًا عدديًا
معاصرًا بوصفه نهائيًا. أما الحد \(H_1\le246\) فيعرض بوصفه حدًا منشورًا موثقًا.

\subsection{نواتج التعلم}

بعد إتمام الفصل ينبغي أن يستطيع القارئ:
\begin{enumerate}
  \item تعريف فجوات الأوليات والكميات \(H_m\).
  \item اشتقاق حد علوي لعدد الأزواج الأولية ذات فرق ثابت من نتائج الفصل الخامس عشر.
  \item استنتاج تقارب مجموع مقلوبات الأزواج الأولية بواسطة الجمع الجزئي.
  \item التمييز بين العدد الأولي والعدد \(P_2\) في مبرهنة تشن.
\personindex{ماينارد، جيمس}
\personindex{تاو، تيرنس}
  \item فهم التحول من GPY إلى Zhang ثم إلى Maynard--Tao.
\personindex{بومبييري، إنريكو}
  \item تفسير دور Bombieri--Vinogradov في طريقة Maynard، ودور التوزيع الأقوى في طريقة Zhang.
\personindex{سيلبرغ، أتله}
  \item التمييز بين غربال سيلبرغ الأحادي والأوزان متعددة الأبعاد.
  \item فهم استمرار عائق التكافؤ رغم نجاح مبرهنات الفجوات المحدودة.
\end{enumerate}

\section{فجوات الأوليات والكميات \(H_m\)}
\symbolindex{فجوات الأوليات والكميات \(H_m\)}

\begin{definition}[فجوات الأوليات]
\label{def:chapter16-prime-gaps}
\resultid{ANT-DEF-16-01}
لتكن \(p_n\) الأولي \(n\)-ي. نكتب
\[
 d_n=p_{n+1}-p_n,
\qquad
 H_m=\liminf_{n\to\infty}(p_{n+m}-p_n),
\qquad m\ge1.
\]
وبخاصة، تكافئ حدسية الأوليات التوأم العبارة \(H_1=2\).
\end{definition}

تقول مبرهنة الأعداد الأولية إن متوسط التباعد قرب \(x\) من رتبة \(\log x\)، لكن
هذا لا يمنع وجود فجوات أصغر كثيرًا من المتوسط أو أكبر كثيرًا منه. السؤال المركزي
هنا هو: هل توجد قيمة مطلقة \(B\) تظهر لا نهائيًا بحيث يوجد أوليان متتاليان فرقُهما
على الأكثر \(B\)؟ أي: هل \(H_1<\infty\)؟

\section{حد علوي للأزواج الأولية}
\symbolindex{حد علوي للأزواج الأولية}

لعدد زوجي ثابت \(h\ne0\)، نعرّف
\[
 \pi_2(x;h)
 =
 \#\{p\le x: p\text{ أولي و }p+h\text{ أولي}\}.
\]

\theoremindex{حد علوي للأزواج الأولية ذات فرق ثابت}
\begin{theorem}[حد علوي للأزواج الأولية ذات فرق ثابت]
\label{thm:chapter16-prime-pair-upper-bound}
\resultid{ANT-THM-16-01}
\provedhere
لكل عدد زوجي ثابت \(h\ne0\)، لدينا
\[
 \pi_2(x;h)\ll_h \frac{x}{(\log x)^2}
 \qquad (x\ge3).
\]
\end{theorem}

\begin{proof}
نختار في مبرهنة الأزواج المنخولة
\ref{thm:chapter15-sifted-pairs} من الفصل الخامس عشر
\[
 z=x^{1/4}.
\]
إذا كان \(p\le x\) و\(p+h\) أوليين، وكان كلاهما أكبر من \(z\)، فإن
\((p(p+h),P(z))=1\). أما الأزواج التي يكون أحد طرفيها أصغر من \(z\) فعددها
\(O_h(z)\). ومن ثم
\[
 \pi_2(x;h)
 \le
 S_h(x,z)+O_h(z).
\]
وبتطبيق المبرهنة \ref{thm:chapter15-sifted-pairs} نحصل على
\[
 S_h(x,z)\ll_h\frac{x}{(\log z)^2}
 \asymp_h\frac{x}{(\log x)^2},
\]
كما أن \(z=x^{1/4}=o(x/(\log x)^2)\)، فتنتج الدعوى.
\end{proof}

هذا حد علوي فقط. لا يثبت أن مجموعة الأزواج غير منتهية، ولا يعطي حدًا سفليًا من
الرتبة نفسها.

\personindex{برون، فيغو}
\section{مجموع برون}
\symbolindex{مجموع برون}

\theoremindex{تقارب مجموع مقلوبات الأزواج الأولية}
\begin{corollary}[تقارب مجموع مقلوبات الأزواج الأولية]
\label{cor:chapter16-brun-sum}
\resultid{ANT-COR-16-01}
\provedhere
لكل عدد زوجي ثابت \(h\ne0\)، تتقارب السلسلة
\[
 \sum_{\substack{p\text{ أولي}\\p+h\text{ أولي}}}\frac1p.
\]
وبخاصة يتقارب مجموع مقلوبات الأوليات التوأم، وهي النتيجة التاريخية المرتبطة
بعمل برون \cite{Brun1919TwinPrimeReciprocals}.
\end{corollary}

\begin{proof}
بالجمع الجزئي، لكل \(X\ge3\)،
\[
 \sum_{\substack{p\le X\\p+h\text{ أولي}}}\frac1p
 =
 \frac{\pi_2(X;h)}{X}
 +
 \int_2^X\frac{\pi_2(t;h)}{t^2}\,dt.
\]
وباستخدام المبرهنة السابقة، يكون الحد الأول \(O_h((\log X)^{-2})\)، بينما
\[
 \int_3^\infty \frac{dt}{t(\log t)^2}<\infty.
\]
إذن تتقارب المجاميع الجزئية.
\end{proof}

المغزى المنهجي مهم: حتى لو كانت الأوليات التوأم غير منتهية، فإنها أندر من أن
تجعل مجموع مقلوباتها يتباعد مثل مجموع مقلوبات جميع الأوليات.

\section{الأعداد شبه الأولية ومبرهنة تشن}
\symbolindex{الأعداد شبه الأولية ومبرهنة تشن}

نقول إن العدد \(n\) من النوع \(P_r\) إذا كان له على الأكثر \(r\) عوامل أولية
مع العد بالتكرار. وبذلك لا يعني \(P_2\) أنه أولي؛ فقد يكون حاصل ضرب أوليين.

\theoremindex{مبرهنة تشن}
\begin{theorem}[مبرهنة تشن]
\label{thm:chapter16-chen}
\resultid{ANT-THM-16-02}
\citedresult{\cite{Chen1973PrimePlusP2}}
كل عدد زوجي كبير بما يكفي يمكن كتابته على الصورة
\[
 N=p+P_2,
\]
حيث \(p\) أولي و\(P_2\) عدد له على الأكثر عاملان أوليان مع العد بالتكرار.
وأثبت تشن أيضًا وجود عدد لا نهائي من الأوليات \(p\) التي يكون
\(p+2\) فيها من النوع \(P_2\).
\end{theorem}

هذه النتيجة تتجاوز غربال سيلبرغ العلوي البسيط؛ فهي تحتاج غربالًا خطيًا وأفكار
تبديل وضبطًا أدق للبواقي. لذلك تُقتبس هنا ولا يُدعى إثباتها كاملًا.

\section{طريقة GPY: فجوات أصغر من المتوسط}
\symbolindex{طريقة GPY: فجوات أصغر من المتوسط}

\theoremindex{Goldston--Pintz--Yıldırım}
\begin{theorem}[Goldston--Pintz--Yıldırım]
\label{thm:chapter16-gpy}
\resultid{ANT-THM-16-03}
\citedresult{\cite{GoldstonPintzYildirim2009PrimesTuplesI}}
لدينا
\[
 \liminf_{n\to\infty}
 \frac{p_{n+1}-p_n}{\log p_n}=0.
\]
\end{theorem}

فكرة GPY هي بناء وزن من نوع سيلبرغ يركز الكتلة على الأعداد \(n\) التي يكون فيها
عدد كبير من الإزاحات \(n+h_i\) قليلة العوامل الأولية الصغيرة، ثم مقارنة مجموع
الوزن بمجموع الوزن المضروب في دالة عد الأوليات على الإزاحات.

لكن النسخة الأصلية من GPY، مع مستوى التوزيع \(1/2\) وحده، لا تعطي مباشرة
\(H_1<\infty\). فإذا تحسن توزيع الأوليات في المتتاليات الحسابية إلى مستوى أكبر
من \(1/2\)، فإن الحجة نفسها تبدأ بإنتاج فجوات محدودة.

\section{اختراق Zhang}
\symbolindex{اختراق Zhang}

\theoremindex{Zhang}
\begin{theorem}[Zhang]
\label{thm:chapter16-zhang}
\resultid{ANT-THM-16-04}
\citedresult{\cite{Zhang2014BoundedGaps}}
أثبت Zhang أن
\[
 H_1<7\times10^7.
\]
وبالتالي توجد لا نهائيًا أزواج من الأوليات المتتالية بفارق محدود مطلق.
\end{theorem}

لم يكن العنصر الحاسم مجرد إعادة تطبيق Bombieri--Vinogradov التقليدية، بل إثبات
توزيع أقوى من المستوى \(1/2\) لعائلة مناسبة من المقامات الملساء، ثم إدخال هذا
التحسن في بنية GPY.

\section{غربال Maynard--Tao متعدد الأبعاد}
\symbolindex{غربال Maynard--Tao متعدد الأبعاد}

التحول التالي لم يكن مجرد تحسين لمستوى التوزيع، بل تغييرًا في بنية الوزن. بدل
وزن أحادي يعتمد أساسًا على حاصل الضرب
\(\prod_i(n+h_i)\)، تُستخدم معاملات متعددة الفهارس
\[
 \lambda_{d_1,\dots,d_k}
\]
وأوزان من الشكل
\[
 w_n=
 \left(
 \sum_{d_i\mid n+h_i}
 \lambda_{d_1,\dots,d_k}
 \right)^2.
\]

بعد اختيار دالة اختبار مناسبة \(F\) على البسيط
\[
 \mathcal R_k=
 \{(t_1,\dots,t_k):t_i\ge0,\ t_1+\cdots+t_k\le1\},
\]
تتحول المسألة إلى مقارنة تكامل رئيسي
\[
 I_k(F)=\int_{\mathcal R_k}F(\mathbf t)^2\,d\mathbf t
\]
بتكاملات هامشية \(J_i(F)\). إذا كانت النسبة التجميعية
\[
 \frac{\sum_{i=1}^kJ_i(F)}{I_k(F)}
\]
كبيرة بما يكفي، أمكن إجبار عدد محدد من الإزاحات على أن تكون أولية في آن واحد.

\theoremindex{Maynard}
\begin{theorem}[Maynard]
\label{thm:chapter16-maynard}
\resultid{ANT-THM-16-05}
\citedresult{\cite{Maynard2015SmallGaps}}
لكل \(m\ge1\)، لدينا
\[
 H_m<\infty.
\]
وفي الحالة \(m=1\)، أعطت الورقة الأصلية الحد
\[
 H_1\le600.
\]
\end{theorem}

الميزة الأساسية أن الطريقة تعمل مع Bombieri--Vinogradov التقليدية، ولا تحتاج
إلى تحسين Zhang الخاص بالمقامات الملساء لإثبات مجرد محدودية الفجوات. وقد طور
Tao الفكرة بصورة مستقلة غير منشورة، بينما المرجع المنشور الحاكم هنا هو عمل
Maynard ثم تعميم Polymath8b \cite{Polymath2014SelbergVariants}.

\section{تحسين Polymath8b}
\symbolindex{تحسين Polymath8b}

\theoremindex{الحد المنشور لـPolymath8b}
\begin{theorem}[الحد المنشور لـPolymath8b]
\label{thm:chapter16-polymath}
\resultid{ANT-THM-16-06}
\citedresult{\cite{Polymath2014SelbergVariants,Polymath2015SelbergVariantsErratum}}
لدينا دون شروط
\[
 H_1\le246.
\]
\personindex{هالبرستام، هيني}
وتحت فرضية Elliott--Halberstam المعممة ثبت الحد
\[
 H_1\le6.
\]
\end{theorem}

يمزج هذا العمل تعميمات لأوزان Maynard متعددة الأبعاد مع تحسينات تحليلية وحسابية.
ولا يُفهم الحد المشروط \(6\) بوصفه برهانًا لحدسية الأوليات التوأم؛ فالانتقال من
\(6\) إلى \(2\) يظل خارج ما تسمح به المعلومات الغربالية البحتة في هذا الإطار.

\section{عائق التكافؤ بعد الفجوات المحدودة}
\symbolindex{عائق التكافؤ بعد الفجوات المحدودة}

لا يعني نجاح Zhang أو Maynard--Tao أن عائق التكافؤ اختفى. الذي تغير هو السؤال:
بدل محاولة إثبات أن نمطًا ثنائيًا محددًا مثل \(n,n+2\) أولي دائمًا على نحو لا
نهائي، نبحث داخل مجموعة كبيرة مقبولة من الإزاحات عن اثنين أو أكثر يقعان معًا في
الأوليات.

يمكن للأوزان متعددة الأبعاد أن تثبت وجود بعض زوج داخل المجموعة، لكنها لا تعزل
بالضرورة زوجًا معينًا. لذلك تبقى حدسية الأوليات التوأم ومسألة \(H_1=2\) مفتوحتين.

\section{خريطة الاعتماد وحدود الادعاء}
\symbolindex{خريطة الاعتماد وحدود الادعاء}

\begin{center}
\begin{tabular}{p{0.23\textwidth}p{0.28\textwidth}p{0.34\textwidth}}
\textbf{النتيجة} & \textbf{المدخل الرئيس} & \textbf{ما لا تستلزمه} \\
\hline
حد الأزواج العلوي & غربال سيلبرغ، الفصل 15 & لا يثبت وجود أزواج لا نهائية \\
تقارب مجموع برون & الحد العلوي والجمع الجزئي & لا يحسم حدسية التوأم \\
تشن & الغربال الخطي والتبديل & لا يجعل \(P_2\) أوليًا \\
GPY & أوزان سيلبرغ وتوزيع المتوسط & لا يعطي وحده فجوة مطلقة مع مستوى \(1/2\) \\
Zhang & GPY وتوزيع أقوى لمقامات ملساء & لا يثبت فجوة \(2\) \\
Maynard--Tao & وزن متعدد الأبعاد وBombieri--Vinogradov & لا يعزل زوجًا محددًا \\
\end{tabular}
\end{center}

\section{خلاصة الفصل}

يعرض هذا الفصل سلسلة من التحولات المنهجية:
\[
 \text{حد علوي غربالي}
 \longrightarrow
 \text{أعداد شبه أولية}
 \longrightarrow
 \text{فجوات صغيرة نسبيًا}
 \longrightarrow
 \text{فجوات محدودة مطلقة}.
\]

أثبتنا داخليًا حد عد الأزواج الأولية وتقارب مجموع برون، واقتبسنا بدقة مبرهنات
تشن وGPY وZhang وMaynard وPolymath8b. وتبقى كل ترقية للحالة مشروطة بالبناء
الكامل، والتدقيق الرياضي، والمراجعة المستقلة.
